\section{Quản lý Scrum và kết quả phát triển qua các giai đoạn}
\label{sec:scrum}
\subsection{Phương pháp và kế hoạch so với kết quả có bằng chứng}
BrewLite được thực hiện theo hướng Agile Scrum: backlog liên kết yêu cầu,
task được xử lý trên branch/PR, kiểm thử và review trước tích hợp. Báo cáo A1
là nguồn mô tả tổ chức trước đây; A2 tiếp nối bằng hai Sprint theo đề.
\path{docs/planning/A2_Implementation_Plan.md} ngày08/11 vẫn ghi PROPOSED/
NOT APPROVED, các mốc là target. Vì vậy không biến lịch mục tiêu thành ngày
Sprint thực tế, không suy PR merged là biên bản Sprint Review.

\begin{longtable}{>{\raggedright\arraybackslash}p{0.22\linewidth}>{\raggedright\arraybackslash}p{0.68\linewidth}}
\caption{Mốc tạo phẩm có bằng chứng GitHub, không thay biên bản Scrum}\label{tab:milestones}\\
\toprule Giai đoạn & Kết quả có nguồn và giới hạn \\ \midrule
\endfirsthead
\toprule Giai đoạn & Kết quả có nguồn và giới hạn (tiếp) \\ \midrule \endhead
A1 kế thừa & Tag A1-final1951dc1, báo cáo A1/contract/ADR giữ provenance.
Vai trò và test A1 không tự trở thành ownership hoặc acceptance A2. \\
Yêu cầu/thiết kế A2 & PR1--9 ngày08--09/11 tạo planning, context/SRS/use case,
activity/sequence/class, architecture/ADR và deployment target. PR merged
chứng minh version history, không phê duyệt mọi gap trong tài liệu. \\
Sprint1 theo Task4 & PR11--15 ngày09--10/11: monorepo, CI, navigation/traceability
và MVP verification report. Report chạy source34cbbc0; không phải E2E production. \\
Hoàn thiện theo Task5 & PR16,18--23,25 ngày10--13/11: proxy/dependency,
Render seed/migration guard, Internal TLS và QueryDR diagnostics. Payment
success claim ở PR26 có phạm vi riêng; không tự đóng SCRUM-114/115. \\
Deployment và delivery & PR24/26/28 promote main; PR27 mô tả CD, ảnh Render
02452f2 và GET bổ sung runtime evidence. PR17 đang chuẩn bị report/release
dossier, Leader chưa review/nghiệm thu. \\
\bottomrule
\end{longtable}

\subsection{Phối hợp, chất lượng và hoạt động còn phải diễn ra}
Git lưu dấu vết thay đổi, PR là nơi reviewer đối chiếu scope và source of
truth; CI chạy install/lint/unit/build để phản hồi sớm. Definition of Done
không chỉ là code có trong branch: cần requirement/test/evidence phù hợp và
quyết định nghiệm thu. Đề còn yêu cầu tạo phẩm các task trước để truy vết task
sau, do đó báo cáo giữ các caveat Approved riêng thay vì tự ghi toàn bộ Done.

Chưa có video end-to-end, biên bản Review/Retro cuối với ngày/người dự/
feedback/action được xác nhận trong các nguồn đã tra. Không tính velocity,
burndown hoặc tỷ lệ hoàn thành khi không có dataset Sprint thực. Dossier có
kịch bản demo và mẫu agenda/biên bản, là Planned. Review cần xem increment,
so với tám AC và ghi quyết định; Retro cần quan sát thực, nguyên nhân, hành
động và người nhận việc sau họp. Nhóm tiếp tục quản lý tài nguyên; tên cá nhân
vận hành chỉ thêm khi được xác nhận.

\section{Đánh giá tổng hợp, hạn chế và hướng phát triển}
\label{sec:evaluation}
\subsection{Kết quả so với mục tiêu đồ án}
Đồ án đã hình thành chuỗi tạo phẩm từ bài toán một cửa hàng đến requirements,
thiết kế, source monorepo, tests/CI và deployment evidence. Tái sử dụng A1
giúp giữ kiến trúc và quy tắc nhất quán; A2 bổ sung khả năng truy vết, chạy
CI, cấu trúc app, migration/seed và tài liệu vận hành. Giá trị của kết quả
nằm ở liên kết nguồn có kiểm chứng, không nằm ở số trang báo cáo.

Audit kế thừa14/11 lúc10:31 UTC củng cố khả năng truy cập HTTP và đọc catalog
tại thời điểm đó, không nâng implementation thành nghiệm thu toàn luồng.
Kết quả cần được đọc cùng lần proxy timeout trước đây ở phần triển khai;
success về sau chưa chứng minh nguyên nhân timeout đã được xử lý.

Các trạng thái được dùng nhất quán: \textbf{Planned} là dự kiến;
\textbf{Implemented} có source/tạo phẩm; \textbf{Verified} có kiểm tra cụ thể
theo SHA/môi trường/thời điểm; \textbf{Approved} có quyết định review đúng phạm
vi; \textbf{Released} có phiên bản/tag và mapping triển khai đã chốt. Một mục
có thể Implemented mà chưa Verified/Approved; CI success không tự nâng thành
Released.

\begin{longtable}{>{\raggedright\arraybackslash}p{0.24\linewidth}>{\raggedright\arraybackslash}p{0.66\linewidth}}
\caption{Đánh giá theo nhóm yêu cầu và lớp bằng chứng}\label{tab:project-assessment}\\
\toprule Nhóm yêu cầu & Kết quả và phần cần hoàn thiện \\ \midrule
\endfirsthead
\toprule Nhóm yêu cầu & Kết quả và phần cần hoàn thiện (tiếp) \\ \midrule \endhead
Catalog/cart/auth & Có UI/API/context và tests. MVP local có report; rounding
preview FE/BE chưa khớp ở giá lẻ. PWA installability chưa có evidence. \\
Order/inventory & Có transaction, state machine, expiry, constraints và test
inventory. Cần runtime integration/concurrency đúng baseline và full UI journey;
không lấy file test làm kết quả test đã chạy. \\
Payment, coupon, loyalty & Có providers/finalizer/idempotency và validation.
Sandbox PAID/SUCCESS/+10 là claim một scenario; status route, IPN/QueryDR E2E,
giao dịch cần đối soát và scope coupon preview còn cần quyết định/evidence. \\
NFR và vận hành & CI lint/unit/build xanh ở head đã commit; HTTP/liveness/catalog
có snapshot. Performance, security measurement, continuous readiness, actual
Internal DB usage, backup/restore và CD chain chưa đủ để nghiệm thu toàn bộ. \\
Delivery/Scrum & Có README/runbook/matrix/report và ảnh đã xử lý; chưa chốt
release/tag, final demo/Review/Retro hoặc Leader acceptance. \\
\bottomrule
\end{longtable}

\subsection{Hạn chế và ưu tiên tiếp theo}
Ưu tiên trước nghiệm thu là chốt scope và giải quyết mismatch requirement--code
có tác động người dùng, thu integration/E2E cho luồng tối thiểu và đối soát
giao dịch còn mở. Về vận hành, cần bind FE production alias/SHA với BE deployment,
xác nhận migration runner và Internal DB config không lộ credential, chốt policy
external access và kế hoạch bảo toàn dữ liệu trước10/12/2026. Sau đó mới chọn
release source/tag và phiên bản rollback đủ điều kiện.

Các cải tiến tương lai như tăng khả dụng scheduler, giảm rủi ro token storage,
PWA hoặc bổ sung ledger phải có requirement/scope/task được duyệt, không tự
triển khai trong phần viết báo cáo. Đo hiệu năng và RPO/RTO cần số liệu thật;
không đặt target tùy ý rồi ghi đã đạt.

Kết quả hiện tại là một \textbf{báo cáo tổng hợp A1 đến A2 có thể review},
kèm tạo phẩm triển khai và giới hạn minh bạch. Báo cáo DRAFT chưa xác nhận
toàn đồ án đã Approved hoặc Released. Quyết định nghiệm thu thuộc Leader/
giảng viên theo đúng nguồn yêu cầu và evidence thực tế.

\subsection{Truy vết và cách đọc phụ lục}
Ma trận chuẩn \path{docs/traceability-matrix.md} nối48 FR/49 NFR với US/UC,
thiết kế, code/API và test inventory. Ví dụ FR-27 nối transaction/order/inventory
với SQL constraints và PostgreSQL concurrency tests; FR-34 nối payment status
với UI/API-03 còn gap; NFR-38 nối CI workflow với run/checks. Đây là mapping
coverage, không phải97 requirements đã PASS. Reviewer phải đọc thêm result/
commit/environment của từng evidence.

Phụ lục giữ ảnh production và bảng readiness tám AC. Nhật ký audit dài, URL/run
IDs, checklist thao tác, kịch bản demo và mẫu biên bản được giữ tại dossier;
không lặp nguyên nhật ký vào nội dung chính. Tài liệu tham khảo ghi nguồn A1,
đề A2, repo và provider; không có tài liệu nguồn nào được upload/chỉnh trực tiếp.
